\section{Introduction}
\label{sec:introduction}

Financial asset returns exhibit occasional movements far too large to be
explained by a Gaussian diffusion. \citet{Merton1976} formalized the now
classical remedy: superimpose a compound jump process on the Brownian
diffusion, so that ordinary days are driven by the diffusion and
extraordinary days by the jumps. Fitting such a model to discretely
sampled data by maximum likelihood is conceptually straightforward ---
write the transition density of an increment, take logs, sum, maximize ---
but almost every step hides a numerical or statistical subtlety:

\begin{itemize}
  \item the transition density involves the \emph{convolution} of the
        Gaussian diffusion with the jump-size law, which has a closed form
        only for a few jump distributions;
  \item the resulting mixture likelihood is \emph{multimodal}, so standard
        gradient-based optimizers can converge to spurious local optima or
        stall at box constraints;
  \item the numerical Hessian at the optimum is frequently ill-conditioned,
        which makes the textbook (Wald) standard errors fragile;
  \item testing whether jumps are present at all puts the null hypothesis
        on the \emph{boundary} of the parameter space while leaving the
        jump-size parameters \emph{unidentified}, so the classical
        likelihood-ratio theory does not apply.
\end{itemize}

The Python library
\code{jump-diffusion-estimation}\footnote{\url{https://github.com/jdospina/jump-diffusion-estimation}}
packages a coherent solution to each of these problems. This working paper
explains \emph{how}: every section states the mathematics first and then
shows the code that implements it, including the design decisions ---
discretization grids, penalization devices, optimizer configurations ---
and the reasoning behind them. Several of those decisions are direct ports
of the author's MSc thesis \citep{Ospina2009} and of the subsequent applied
work in \citet{Ardia2011}; where that is the case we say so explicitly, so
the reader can trace each formula to its origin.

\subsection{The library at a glance}
\label{sec:architecture}

The package is organized in four cooperating layers, each mapped to a
Python subpackage:

\begin{center}
\small
\begin{tabular}{@{}llp{5.4cm}@{}}
\toprule
Layer & Subpackage & Main classes \\
\midrule
Model \& simulation & \code{models}, \code{simulation} &
  \code{JumpDiffusionModel},\newline \code{JumpDiffusionSimulator} \\
Jump distributions & \code{distributions} &
  \code{JumpDistribution} + five concrete families \\
Estimation \& inference & \code{estimation} &
  \code{JumpDiffusionEstimator} \\
Validation \& comparison & \code{validation} &
  \code{ValidationExperiment},\newline \code{JumpDistributionComparison} \\
\bottomrule
\end{tabular}
\end{center}

The dependency direction is strictly downward in the table:
\code{JumpDiffusionModel} owns the transition density and path simulation,
delegating everything jump-specific to a \code{JumpDistribution} object;
\code{JumpDiffusionEstimator} evaluates that model's likelihood at trial
parameter vectors and wraps the optimizers and the inference machinery;
the validation layer repeatedly drives the estimator on simulated or real
data. A user who wants a new jump-size family implements a single class
with the distribution's density --- everything else (likelihood, FFT
convolution, sampling, estimation, inference, comparison) works unchanged.
That extension contract is the subject of Section~\ref{sec:distributions}.

\subsection{Notation}
\label{sec:notation}

Throughout the paper, $X_t$ denotes the (log-price) process, $\mu$ the
drift, $\sigma > 0$ the diffusion volatility, $W_t$ a standard Brownian
motion, $N_t$ a Poisson counting process with intensity $\lambda$, and
$J$ a generic jump size with density $f_J(\cdot\,;\theta_J)$ and parameter
vector $\theta_J$. Observations are $n$ increments
$\Delta x_i = x_{t_i} - x_{t_{i-1}}$ on a regular grid with spacing
$\Delta t$; in the code, \code{data} is the array of increments and
\code{dt} is $\Delta t$. The Bernoulli jump probability per interval is
$p$ (\code{jump\_prob} in the code), which plays the role of
$\lambda \Delta t$. The full parameter vector is
\[
  \thetavec \;=\; (\mu,\ \sigma,\ p,\ \theta_J),
\]
and the code orders it exactly this way:
\begin{center}
\code{("mu", "sigma", "jump\_prob", *jump\_distribution.param\_names)}.
\end{center}
We write $\phi(x; m, s^2)$ for the normal density with mean $m$ and
variance $s^2$, $\Phi$ for the standard normal CDF, $*$ for convolution,
and $\loglik(\thetavec)$ for the log-likelihood. Code identifiers are set
in \code{typewriter} font; all listings are excerpts of the actual library
source (docstrings and defensive branches are sometimes elided for
readability, never altered).

\subsection{How to read this paper}

Sections follow the pipeline of an actual analysis.
Section~\ref{sec:model} derives the model and its Bernoulli-approximated
transition density --- the single formula the whole library revolves
around. Section~\ref{sec:distributions} presents the five jump-size
families and their closed-form Gaussian convolutions where they exist.
Section~\ref{sec:fft} derives the FFT convolution scheme used when they
do not. Section~\ref{sec:estimation} covers maximum likelihood and the
two optimizers. Section~\ref{sec:inference} derives the three
uncertainty-quantification routes, Section~\ref{sec:jumptest} the
bootstrap test for the presence of jumps, and Section~\ref{sec:comparison}
the goodness-of-fit comparison and Monte Carlo validation tools.
Section~\ref{sec:closing} closes with the research agenda this library
supports.
